% Blueprint chapter for VerifiedReserving/ODPStochastic.lean and its finite witness.

\chapter{The stochastic ODP layer}

England and Verrall describe independent incremental claims with
\[
  \mathbb{E}[X_{i,k}]=m_{i,k}, \qquad
  \operatorname{Var}(X_{i,k})=\varphi m_{i,k}, \qquad
  \log m_{i,k}=c+\alpha_i+\beta_k,
\]
with the corner constraints $\alpha_0=\beta_0=0$
\cite[Section 2.2, equations (2.5)--(2.7)]{EnglandVerrall1999}.
The same moment and link equations appear in
\cite[Section 2, equations (2.1)--(2.3)]{England2002}.
This chapter separates exact consequences of those assumptions from the first-order GLM
prediction-error formula and the residual bootstrap algorithm.

\section{Model and conditional moments}

\begin{definition}[Finite stochastic ODP model]\label{def:odpStochasticModel}\lean{VerifiedReserving.ODPCell, VerifiedReserving.odpCellRv, VerifiedReserving.ODPModel, VerifiedReserving.odpCellSigma, VerifiedReserving.odpObservedRv, VerifiedReserving.odpObservedSigma}\leanok
For a finite square of cells, every incremental claim is measurable and square-integrable,
has mean $m_{i,k}$ and variance $\varphi m_{i,k}$, and the family of incremental claims is
mutually independent. The observed-data sigma-algebra is generated by any selected finite
set of cells.
\end{definition}

\begin{lemma}[Generated sigma-algebras and independence]\label{lem:odpObservedIndep}\lean{VerifiedReserving.odpCellSigma_le, VerifiedReserving.odpObservedSigma_le, VerifiedReserving.odpIndep_cell_observed}\leanok
\uses{def:odpStochasticModel}
The sigma-algebras generated by one cell and by a finite vector of cells are sub-sigma-algebras
of the ambient one. A cell outside the selected vector is independent of that vector.
\end{lemma}
\begin{proof}\leanok Measurability gives the two inclusions. Mutual independence, restricted to
the disjoint singleton and observed index sets, gives the final statement. \end{proof}

\begin{theorem}[Exact conditional ODP moments]\label{thm:odpConditionalMoments}\lean{VerifiedReserving.condExp_odpCell_given_observed, VerifiedReserving.condExp_sq_odpCell_sub_given_observed}\leanok
\uses{lem:odpObservedIndep}
If $(i,k)$ is not among the observed cells, then
\[
  \mathbb{E}[X_{i,k}\mid\mathcal{O}]=m_{i,k}, \qquad
  \mathbb{E}[(X_{i,k}-m_{i,k})^2\mid\mathcal{O}]=\varphi m_{i,k}
\]
almost surely. These are exact conditional-moment identities, not GLM approximations.
\end{theorem}
\begin{proof}\leanok Independence makes each conditional expectation equal to its integral;
substitute the two fields of the ODP model. \end{proof}

\section{Prediction error}

\begin{definition}[Aggregate and mean squared prediction error]\label{def:odpMsepExact}\lean{VerifiedReserving.odpAggregate, VerifiedReserving.meanSquaredPredictionError}\leanok
For a finite set $F$ of future cells, let $Y_F=\sum_{p\in F}X_p$. For a predictor
$\widehat Y$, define $\operatorname{MSEP}(Y_F,\widehat Y)
=\mathbb{E}[(Y_F-\widehat Y)^2]$.
\end{definition}

\begin{theorem}[Aggregate mean and process variance]\label{thm:odpAggregateMoments}\lean{VerifiedReserving.integral_odpAggregate, VerifiedReserving.variance_odpAggregate}\leanok
\uses{def:odpStochasticModel, def:odpMsepExact}
For every finite $F$,
\[
  \mathbb{E}[Y_F]=\sum_{p\in F}m_p,
  \qquad
  \operatorname{Var}(Y_F)=\varphi\sum_{p\in F}m_p.
\]
\end{theorem}
\begin{proof}\leanok Integrate the finite sum and add the variances of the independent
square-integrable cells. \end{proof}

\begin{theorem}[Exact process-plus-estimation decomposition]\label{thm:odpMsepExact}\lean{VerifiedReserving.meanSquaredPredictionError_eq_variance_add, VerifiedReserving.odpAggregateMsep_eq}\leanok
\uses{def:odpMsepExact, thm:odpAggregateMoments}
If a square-integrable predictor has the same mean as its square-integrable target and is
independent of it, then
\[
  \mathbb{E}[(Y-\widehat Y)^2]=\operatorname{Var}(Y)+\operatorname{Var}(\widehat Y).
\]
For an ODP aggregate this is
$\varphi\sum_{p\in F}m_p+\operatorname{Var}(\widehat Y)$.
\end{theorem}
\begin{proof}\leanok The difference has mean zero. Expand its variance and use independence to
set the covariance to zero, then substitute Theorem~\ref{thm:odpAggregateMoments}. \end{proof}

England and Verrall use the delta method for the estimation component: since
$\widehat m=\exp(\widehat\eta)$,
$\operatorname{Var}(\widehat m)\mathrel{\dot=}\widehat m^2
\operatorname{Var}(\widehat\eta)$. Their single-cell and aggregate formulas are equations
(3.3)--(3.5) of \cite{EnglandVerrall1999} and equations (7.7)--(7.9) of
\cite{EnglandVerrall2002}. The dotted equality is not promoted to a theorem.

\begin{definition}[Delta-method GLM MSEP formulas]\label{def:odpGlmMsep}\lean{VerifiedReserving.glmCellMsep, VerifiedReserving.glmAggregateEstimationVariance, VerifiedReserving.glmAggregateMsep}\leanok
The single-cell formula is
$\varphi\widehat m+\widehat m^2\operatorname{Var}(\widehat\eta)$.
For a finite aggregate, the estimation term is the covariance quadratic form
$\sum_{p,q}\widehat m_p\widehat m_q\operatorname{Cov}(\widehat\eta_p,\widehat\eta_q)$,
added to $\varphi\sum_p\widehat m_p$.
These are definitions of the source approximation.
\end{definition}

\begin{lemma}[Singleton bridge]\label{lem:odpGlmSingleton}\lean{VerifiedReserving.glmAggregateMsep_singleton}\leanok
\uses{def:odpGlmMsep}
On a singleton future set, the aggregate definition is the single-cell formula, with the
diagonal covariance used as $\operatorname{Var}(\widehat\eta)$.
\end{lemma}
\begin{proof}\leanok Evaluate both finite sums on a singleton. \end{proof}

\section{Log-link and corner constraints}

\begin{definition}[Log-link and corner effects]\label{def:odpLogLink}\lean{VerifiedReserving.logLinkFit, VerifiedReserving.cornerRow, VerifiedReserving.cornerColumn}\leanok
The log-link fit is $\exp(c+\alpha_i+\beta_k)$. For positive multiplicative parameters
$a_i,b_k$, set
$\alpha_i=\log a_i-\log a_0$ and $\beta_k=\log b_k-\log b_0$.
\end{definition}

\begin{theorem}[Log-link and multiplicative equivalence]\label{thm:odpLogLink}\lean{VerifiedReserving.logLinkFit_eq_multFit, VerifiedReserving.cornerRow_zero, VerifiedReserving.cornerColumn_zero, VerifiedReserving.logLinkFit_corner_eq_multFit}\leanok
\uses{def:odpLogLink, def:mult}
Every log-link fit is multiplicative. Conversely, every positive multiplicative fit has the
corner-constrained representation above, with $c=\log a_0+\log b_0$ and
$\alpha_0=\beta_0=0$.
\end{theorem}
\begin{proof}\leanok Use $\exp(x+y)=\exp(x)\exp(y)$ in one direction and
$\exp(\log x)=x$ for positive $x$ in the other. \end{proof}

\section{Residual bootstrap definitions}

England and Verrall's unscaled Pearson residual, scale estimate, and inversion are equations
(4.2)--(4.4) of \cite{EnglandVerrall1999}. England's later two-stage procedure first resamples
adjusted residuals and refits the mean, then simulates future process outcomes conditional on
that mean \cite[Sections 2--3]{England2002}.

\begin{definition}[Pearson residual operations]\label{def:odpBootstrapResidual}\lean{VerifiedReserving.pearsonResidual, VerifiedReserving.pearsonScale, VerifiedReserving.adjustedPearsonResidual, VerifiedReserving.bootstrapIncremental}\leanok
For an observation $x$ and fitted mean $m$, the unscaled Pearson residual is
$r=(x-m)/\sqrt m$. The Pearson scale is $\sum r^2/(N-p)$, the degrees-of-freedom
adjustment multiplies a residual by $\sqrt{N/(N-p)}$, and inversion gives the pseudo-observation
$x^*=r^*\sqrt m+m$. Lean's total division leaves these definitions meaningful at zero; the
inversion theorems below assume $m>0$.
\end{definition}

\begin{lemma}[Pearson inversion]\label{lem:odpBootstrapInverse}\lean{VerifiedReserving.pearsonResidual_bootstrapIncremental, VerifiedReserving.bootstrapIncremental_pearsonResidual}\leanok
\uses{def:odpBootstrapResidual}
For $m>0$, taking the residual of $r\sqrt m+m$ returns $r$, and rebuilding an observation
from its Pearson residual returns that observation.
\end{lemma}
\begin{proof}\leanok Cancel the nonzero factor $\sqrt m$. \end{proof}

\begin{definition}[England-Verrall bootstrap and predictive reserve]\label{def:odpBootstrap}\lean{VerifiedReserving.englandVerrallBootstrapMean, VerifiedReserving.englandVerrallPredictiveReserve}\leanok
One bootstrap refit resamples supplied residuals with replacement, inverts them against the
fitted past means, and applies a supplied refitting map to obtain future means. The predictive
reserve sums supplied process draws conditional on those bootstrap means.
The process simulator is an input: the ODP moment equations do not determine a unique
distribution. These declarations define the algorithm and make no coverage or consistency claim.
\end{definition}

\section{Finite non-vacuity witness}

\begin{theorem}[A nondegenerate ODP model exists]\label{thm:odpWitness}\lean{VerifiedReserving.ODPWitness.Omega, VerifiedReserving.ODPWitness.mu, VerifiedReserving.ODPWitness.mu_singleton, VerifiedReserving.ODPWitness.X, VerifiedReserving.ODPWitness.m, VerifiedReserving.ODPWitness.phi, VerifiedReserving.ODPWitness.integral_X, VerifiedReserving.ODPWitness.variance_X, VerifiedReserving.ODPWitness.model, VerifiedReserving.ODPWitness.p00, VerifiedReserving.ODPWitness.nontrivial, VerifiedReserving.ODPWitness.conditionalMean, VerifiedReserving.ODPWitness.conditionalVariance, VerifiedReserving.ODPWitness.exists_nondegenerate_odp_model}\leanok
\uses{thm:odpConditionalMoments}
On a two-point uniform probability space, let the only modeled cell be $0$ on one outcome and
$2$ on the other. It has mean and variance one, is genuinely random, satisfies the ODP model
with $m=\varphi=1$, and instantiates both exact conditional-moment theorems.
\end{theorem}
\begin{proof}\leanok Evaluate the two finite sums; the single-cell independence condition is
automatic. \end{proof}
